% !TeX program = xelatex
% !TeX encoding = UTF-8
% 在 out 目录中运行：xelatex -interaction=nonstopmode -halt-on-error ../main.tex

\PassOptionsToPackage{no-math}{fontspec}
\documentclass[
  12pt,
  UTF8,
  a4paper,
  fontset=fandol,
  oneside,
  openany
]{ctexrep}

\AddToHook{package/arydshln/before}{
  \RequirePackage{longtable}
  \RequirePackage[table]{xcolor}
}

\RequirePackage{amsmath}
\RequirePackage{array}
\RequirePackage{arydshln}
\RequirePackage{booktabs}
\RequirePackage{caption}
\RequirePackage{circuitikz}
\RequirePackage{float}
\RequirePackage[
  bottom=2.54cm,
  left=3.18cm,
  right=3.18cm,
  top=2.54cm
]{geometry}
\RequirePackage[hidelinks,unicode]{hyperref}
\RequirePackage{longtable}
\RequirePackage{subcaption}
\RequirePackage{unicode-math}
\RequirePackage[table]{xcolor}

\usepackage{amsmath}
\usepackage{array}
\usepackage{arydshln}
\usepackage{booktabs}
\usepackage{caption}
\usepackage{circuitikz}
\usepackage{float}
\usepackage[
  bottom=2.54cm,
  left=3.18cm,
  right=3.18cm,
  top=2.54cm
]{geometry}
\usepackage[hidelinks,unicode]{hyperref}
\usepackage{longtable}
\usepackage{subcaption}
\usepackage{unicode-math}
\usepackage[table]{xcolor}

\definecolor{reviewgray}{gray}{0.42}
\DeclareCaptionFont{wuhao}{\zihao{5}}
\DeclareCaptionLabelFormat{figurelabel}{图#2}
\DeclareCaptionLabelFormat{tablelabel}{表#2}
\DeclareCaptionLabelSeparator{figurespace}{\hspace{1em}}
\captionsetup[figure]{
  font=wuhao,
  justification=centering,
  labelformat=figurelabel,
  labelsep=figurespace
}
\captionsetup[subfigure]{
  font=wuhao,
  justification=centering,
  labelformat=parens,
  labelsep=space
}
\captionsetup[table]{
  font=wuhao,
  justification=centering,
  labelformat=tablelabel,
  labelsep=figurespace
}
\ctikzset{
  bipoles/thickness=1,
  inductor=american,
  quadpoles/thickness=1,
  tripoles/thickness=1
}
\hypersetup{pdftitle={电力系统分析笔记}}
\linespread{1.3}
\pagestyle{plain}
\raggedbottom{}
\setcounter{secnumdepth}{5}
\setcounter{tocdepth}{3}
\setlength{\parindent}{2em}
\setlength{\parskip}{0pt}
\setmathfont{XITS Math}
\setmathfont[version=bold]{XITS Math}

\newcounter{example}[chapter]
\newcounter{examplesubheading}[example]
\newfontfamily\circlenumfont{Times New Roman}
\newcommand{\exampletitle}[2][]{%
  \refstepcounter{example}%
  \par\addvspace{0.15cm}%
  \noindent{\heiti\zihao{-4}\bfseries
    ※\hspace{0.3em}例\theexample\hspace{1em}#2}%
  \if\relax\detokenize{#1}\relax\else\label{#1}\fi
  \par\nobreak{}
}
\newcommand{\examplesubheading}[1]{%
  \stepcounter{examplesubheading}%
  \ifnum\value{examplesubheading}=1
    \begingroup
      \ctexset{subsubsection={beforeskip=0pt}}%
      \subsubsection*{%
        \timescircled{\arabic{examplesubheading}}%
        \hspace{0.5em}#1%
      }%
    \endgroup
  \else
    \subsubsection*{%
      \timescircled{\arabic{examplesubheading}}%
      \hspace{0.5em}#1%
    }%
  \fi
}
\newcommand{\notescircuitfont}{\fontsize{9}{10.8}\selectfont}
\newenvironment{notescircuit}[1][]{%
  \begingroup
  \notescircuitfont
  \ctikzset{
    bipoles/length=2.276em,
    transistors/scale=1.5
  }
  \begin{circuitikz}[
    every node/.append style={font=\notescircuitfont},
    font=\notescircuitfont,
    line width=0.8pt,
    transform shape=false,
    #1
  ]
}{%
  \end{circuitikz}
  \endgroup
}
\newcommand{\reviewnote}[1]{%
  \par\smallskip
  {\color{reviewgray}\small\noindent【校核建议：#1】\par}
  \smallskip
}
\newcommand{\figureplaceholder}[2]{%
  \par\smallskip
  {\color{reviewgray}\small\noindent【图像预留 #1：#2】\par}
  \vspace{0.5cm}
}
\newcommand{\timescircled}[1]{%
  \raisebox{0.5pt}{%
    \tikz[
      baseline=(circlednumber.base),
      scale=0.82,
      transform shape
    ]{%
      \node[
        circle,
        draw,
        font=\circlenumfont\mdseries,
        inner sep=0pt,
        line width=0.4pt,
        minimum size=1.15em
      ] (circlednumber) {#1};%
    }%
  }%
}
\renewcommand{\thechapter}{\arabic{chapter}}
\renewcommand{\theequation}{\arabic{chapter}-\arabic{equation}}
\renewcommand{\theexample}{\thechapter-\arabic{example}}
\renewcommand{\thefigure}{\arabic{chapter}-\arabic{figure}}
\renewcommand{\thetable}{\arabic{chapter}-\arabic{table}}
\renewcommand{\thesection}{\thechapter.\arabic{section}}
\renewcommand{\thesubsection}{\thesection.\arabic{subsection}}
\renewcommand{\thesubsubsection}{\thesubsection.\arabic{subsubsection}}
\renewcommand{\theparagraph}{\thesubsubsection.\arabic{paragraph}}
\renewcommand{\thesubparagraph}{\theparagraph.\arabic{subparagraph}}

\ctexset{
  chapter={
    afterindent=true,
    aftername=\quad,
    afterskip=1.5ex,
    beforeskip=0pt,
    format=\heiti\centering\zihao{2}\bfseries,
    name={第,章},
    number=\chinese{chapter}
  },
  section={
    afterindent=true,
    aftername=\hspace{0.5em},
    afterskip=1ex,
    beforeskip=1ex,
    format=\heiti\raggedright\zihao{3}\bfseries,
    name={\S,},
    number=\thesection
  },
  subsection={
    afterindent=true,
    aftername=\hspace{0.5em},
    afterskip=0ex,
    beforeskip=0.3cm,
    format=\heiti\raggedright\zihao{4}\bfseries,
    name={,},
    number=\thesubsection
  },
  subsubsection={
    afterindent=true,
    aftername=\hspace{0.5em},
    afterskip=0ex,
    beforeskip=0.15cm,
    format=\songti\raggedright\zihao{-4}\bfseries,
    name={,},
    number=\thesubsubsection
  },
  paragraph={
    afterindent=true,
    aftername=\hspace{0.5em},
    afterskip=0ex,
    beforeskip=0.5ex,
    format=\songti\raggedright\zihao{-4},
    indent=0pt,
    name={,},
    number=\theparagraph,
    runin=false
  },
  subparagraph={
    afterindent=true,
    aftername=\hspace{0.5em},
    afterskip=0ex,
    beforeskip=0.5ex,
    format=\songti\raggedright\zihao{-4},
    indent=0pt,
    name={,},
    number=\thesubparagraph,
    runin=false
  },
}

\IfFontExistsTF{SimHei}{
  \setCJKfamilyfont{zhhei}[AutoFakeBold=2]{SimHei}
  \setCJKsansfont[AutoFakeBold=2]{SimHei}
}{}
\IfFontExistsTF{SimSun}{
  \setCJKfamilyfont{zhsong}[AutoFakeBold=2]{SimSun}
  \setCJKmainfont[AutoFakeBold=2]{SimSun}
}{}
\IfFontExistsTF{Times New Roman}{
  \setmainfont{Times New Roman}
}{
  \setmainfont{TeX Gyre Termes}
}{}

\ExplSyntaxOn
\prop_new:N \g_notes_block_owner_prop
\prop_new:N \g_notes_block_declared_prop
\prop_new:N \g_notes_block_defined_prop
\prop_new:N \g_notes_block_used_prop

\NewDocumentCommand{\DeclareContentBlock}{m m}{
  \clist_set:Nn \l_tmpa_clist {#2}
  \tl_set:Nx \l_tmpa_tl {\clist_item:Nn \l_tmpa_clist {1}}
  \clist_map_inline:Nn \l_tmpa_clist {
    \prop_if_in:NnTF \g_notes_block_owner_prop {#1/##1}
      {\PackageError{notes}{Duplicate~#1~declaration:~##1}{Use~a~unique~label~in~the~header.}}
      {\prop_gput:NnV \g_notes_block_owner_prop {#1/##1} \l_tmpa_tl}
  }
  \exp_args:NNx \prop_gput:Nnn \g_notes_block_declared_prop
    {#1/\tl_use:N \l_tmpa_tl} {true}
  \tl_new:c {g_notes_block_#1_\tl_use:N \l_tmpa_tl _tl}
}

\NewDocumentCommand{\DefineContentBlock}{m m +m}{
  \prop_if_in:NnTF \g_notes_block_declared_prop {#1/#2} {
    \prop_if_in:NnTF \g_notes_block_defined_prop {#1/#2}
      {\PackageError{notes}{Duplicate~#1~definition:~#2}{Define~each~block~only~once.}}
      {
        \tl_gset:cn {g_notes_block_#1_#2_tl} {#3}
        \prop_gput:Nnn \g_notes_block_defined_prop {#1/#2} {true}
      }
  }{\PackageError{notes}{Undeclared~#1~block:~#2}{Declare~the~first~label~of~the~block~in~the~header.}}
}

\cs_new_protected:Npn \notes_insert_block:nn #1#2 {
  \prop_if_in:NnTF \g_notes_block_defined_prop {#1/#2} {
    \prop_if_in:NnTF \g_notes_block_used_prop {#1/#2}
      {\PackageError{notes}{Repeated~#1~insertion:~#2}{Use~ref~or~eqref~to~refer~to~an~existing~block.}}
      {
        \prop_gput:Nnn \g_notes_block_used_prop {#1/#2} {true}
        \tl_use:c {g_notes_block_#1_#2_tl}
      }
  }{\PackageError{notes}{Undefined~#1~block:~#2}{Check~the~corresponding~source~definition.}}
}

\NewDocumentCommand{\InsertContentBlock}{m m}{
  \prop_get:NnNTF \g_notes_block_owner_prop {#1/#2} \l_tmpa_tl
    {\exp_args:NnV \notes_insert_block:nn {#1} \l_tmpa_tl}
    {\PackageError{notes}{Unknown~#1~label:~#2}{Check~the~corresponding~header.}}
}

\cs_new_protected:Npn \notes_check_definitions: {
  \prop_map_inline:Nn \g_notes_block_declared_prop {
    \prop_if_in:NnF \g_notes_block_defined_prop {##1}
      {\PackageError{notes}{Missing~block~definition:~##1}{Load~all~declared~source~files~before~the~document.}}
  }
}
\AtBeginDocument{\notes_check_definitions:}
\ExplSyntaxOff

\tikzset{
  notesline/.style={line width=0.65pt},
  notesbus/.style={line width=1pt},
  notesflow/.style={-{Latex[length=1.8mm]}, line width=0.65pt},
  pics/notesgenerator/.style={code={
    \draw[notesline] (0,0) circle[radius=0.38];
    \node at (0,0.10) {$G$};
    \node at (0,-0.14) {$\sim$};
  }},
  pics/notesmotor/.style={code={
    \draw[notesline] (0,0) circle[radius=0.38];
    \node at (0,0.10) {$M$};
    \node at (0,-0.14) {$\sim$};
  }},
  pics/notestwowinding/.style={code={
    \draw[notesline] (-0.22,0) circle[radius=0.32];
    \draw[notesline] (0.22,0) circle[radius=0.32];
  }},
  pics/notesthreewinding/.style={code={
    \draw[notesline] (-0.22,0) circle[radius=0.32];
    \draw[notesline] (0.22,0) circle[radius=0.32];
    \draw[notesline] (0,-0.36) circle[radius=0.32];
  }},
  pics/notesbusbar/.style={code={
    \draw[notesbus] (0,-0.36) -- (0,0.36);
  }}
}

\NewDocumentCommand{\DeclareEquation}{m}{\DeclareContentBlock{equation}{#1}}
\NewDocumentCommand{\DefineEquation}{m +m}{\DefineContentBlock{equation}{#1}{#2}}
\NewDocumentCommand{\InsertEquation}{m}{\InsertContentBlock{equation}{#1}}

% 公式声明预留；统一标签后再注册。
% \DeclareEquation{equ0101}

% 第一章公式定义预留。
% \DefineEquation{equ0101}{
% }

\NewDocumentCommand{\DeclareFigure}{m}{\DeclareContentBlock{figure}{#1}}
\NewDocumentCommand{\DefineFigure}{m +m}{\DefineContentBlock{figure}{#1}{#2}}
\NewDocumentCommand{\InsertFigure}{m}{\InsertContentBlock{figure}{#1}}

% 图像声明预留；统一标签后再注册。
% \DeclareFigure{pic0101}

% 第一章图像定义预留。
% \DefineFigure{pic0101}{
% }

\NewDocumentCommand{\DeclareTable}{m}{\DeclareContentBlock{table}{#1}}
\NewDocumentCommand{\DefineTable}{m +m}{\DefineContentBlock{table}{#1}{#2}}
\NewDocumentCommand{\InsertTable}{m}{\InsertContentBlock{table}{#1}}

% 表格声明预留；统一标签后再注册。
% \DeclareTable{tab0101}

% 第一章表格定义预留。
% \DefineTable{tab0101}{
% }






% ---------- 正文 ----------

\begin{document}

% 第一章正文预留；新材料先在 main.tex 整理，授权分文件后再迁入。
% \chapter{待定}
% \section{待定}


% 手稿 1：MsgID=803879089740954952，文件修改时间 2026-10-07 19:24:06.303。
% 四张图片均无可用的 EXIF 拍摄时间，按文件修改时间排序并核对上下文。
\chapter{电力系统基本概念}

\section{电力系统的组成}

\subsection{电能的优势}

\reviewnote{原稿仅注明“书 P1”，未列出具体内容；待依据所用教材补录。}

\subsection{电力系统的定义}

\reviewnote{原稿仅注明“书 P1”，未写出定义；待依据所用教材补录。}

\subsection{电力系统简图与主要组成}

电力系统的发电、输送、分配和消费环节，可以用单线图表示。

% 原图 1-A。
\begin{figure}[H]
  \centering
  \begin{tikzpicture}[every node/.style={font=\small}, notesline]
    \draw (0.10,-0.28) -- (0.75,-0.17) -- (0.75,0.17)
    -- (0.10,0.28) -- cycle;
    \draw[dashed] (0.75,0) -- (1.27,0);
    \pic at (1.65,0) {notesgenerator};
    \pic at (3.50,0) {notestwowinding};
    \pic at (8.00,0) {notestwowinding};
    \draw (2.03,0) -- (2.96,0);
    \draw (4.04,0) -- (7.46,0);
    \draw (8.54,0) -- (12.20,0);
    \draw[notesflow] (12.20,0) -- (12.20,-0.45);
    \foreach \x in {2.45,4.50,7.00,9.00,11.50}
    \pic at (\x,0) {notesbusbar};
    \node at (0.43,-0.85) {汽轮机};
    \node at (1.65,-0.85) {发电机};
    \node at (3.50,-0.85) {升压变压器};
    \node at (5.75,-0.85) {输电线路};
    \node at (8.00,-0.85) {降压变压器};
    \node at (10.20,-0.85) {配电线路};
    \node at (12.20,-0.85) {负荷};
    \draw (0.10,0.95) -- (0.10,1.08) -- (2.45,1.08) -- (2.45,0.95);
    \draw (2.65,0.95) -- (2.65,1.08) -- (7.00,1.08) -- (7.00,0.95);
    \draw (7.20,0.95) -- (7.20,1.08) -- (11.50,1.08) -- (11.50,0.95);
    \draw (11.70,0.95) -- (11.70,1.08) -- (12.65,1.08) -- (12.65,0.95);
    \node at (1.28,1.40) {电源};
    \node at (4.83,1.40) {输送};
    \node at (9.35,1.40) {分配};
    \node at (12.18,1.40) {消费};
  \end{tikzpicture}
  \caption{电力系统的基本组成与单线示意}
\end{figure}

\reviewnote{本图按原稿概括电源、输送、分配和消费环节，各环节上方的括线用于说明流程，不能作为输电与配电的严格电压分界。汽轮机与发电机之间的虚线表示机械联接；电气部分的单线代表三相系统。更复杂的接线及包含关系仍需核对教材 P12。}

\begin{itemize}
  \item 发电机：发电设备，将其他形式的能量转化为电能。
  \item 变压器：原稿记为降低有功（热）损耗、减小电压降、提升稳定性。
  \item 线路：用于电能的输送和分配。
  \item 负荷：用电设备，将电能转化为其他形式的能量；原稿另记“大部分转化为内能（发热）”。
\end{itemize}

\reviewnote{变压器的作用应与升压输电的效果区分。给定传输功率及功率因数时，提高输电电压可减小电流和线路有功损耗；变压器本身仍有损耗。“减小电压降、提升稳定性”也需结合线路参数和运行条件说明，不宜作为所有变压器的无条件结论。}

\reviewnote{负荷的有用能量输出与最终耗散应区分。“大部分发热”若指最终转化为内能，需要明确这一含义；不能将电动机的机械输出等一概视为用电设备的损耗。}

\subsection{一次系统与二次系统}

上述单线图中的发电、变压、输配电及用电部分构成一次系统；用于监测、保护、控制和调节的部分构成二次系统。

原稿指出，二次系统可通过互感线圈，将一次系统中的大电压或大电流转换为便于测量的小电压或小电流，其命名与变压器有相似之处，参见教材 P3。

\reviewnote{这里应区分电压互感器和电流互感器。“一次系统/二次系统”是按设备功能分类，与变压器的“一次侧/二次侧”不是同一概念；建议后续补充互感器在测量及保护中的作用。}

\subsection{直流输电系统}

\reviewnote{原稿仅注明“书 P2”，未展开直流输电系统的结构与特点；暂不补造内容。}

\section{电力系统的运行特点与基本要求}

\subsection{运行特点}

原稿列出的运行特点如下，参见教材 P4。

\begin{enumerate}
  \item 电能不能大量存储，发电、输电、配电和用电必须同时完成。
  \item 输送快、传播快。
  \item 与生产、生活密切相关。
\end{enumerate}

\reviewnote{“不能大量存储”应限定在传统电力系统的常规运行描述中。电池、抽水蓄能等可通过能量转换实现储存；更准确的表述是系统仍需保持实时功率平衡，储能可参与这一平衡。参见\href{https://www.energy.gov/energy-storage}{美国能源部储能说明}。}

\subsection{基本要求}

\begin{itemize}
  \item 具有备用及储能手段，保证电能质量，主要涉及电压、频率和波形。
  \item 监测、控制和保护应快速、准确，并保证系统稳定运行。
  \item 保证供电可靠性；从技术和经济角度，负荷分为三级，参见教材 P4。
  \item 保证经济性，降低煤耗率、水耗率及网损率。
  \item 满足低碳、环保要求，关注新能源的发展。
\end{itemize}

\textbf{疑问：}分布式还是集中式？

\reviewnote{原稿在低碳、环保要求旁提出上述问题，未给出结论。分布式与集中式的选择需结合能源分布、负荷位置及电网条件，后续再展开。}

% 手稿 2：MsgID=6375615577849184564，文件修改时间 2026-10-07 19:24:09.828。
原稿指出，新能源具有间歇性、波动性等不可控特点，对电力系统提出了挑战。

\reviewnote{“间歇性”之后另有一项“地……性”的字迹不清，疑似“地域性”，待核对。上述描述主要适用于风能、太阳能等能源的可用出力，不宜将所有新能源的出力和控制能力一概归为不可控。}

\section{电力系统的额定电压}

本节原稿参考教材 P2-P4，并注明“均为线电压”。应区分系统额定电压与设备额定电压。

\subsection{电压等级}

原稿按低压、中压、高压、超高压和特高压列出以下五组电压等级。

\begin{enumerate}
  \item 低压：三相 $380~\mathrm{V}$，单相 $220~\mathrm{V}$。
  \item 中压：$3~\mathrm{kV}$、$6~\mathrm{kV}$、$10~\mathrm{kV}$、$20~\mathrm{kV}$、$35~\mathrm{kV}$、$66~\mathrm{kV}$。
  \item 高压：$110~\mathrm{kV}$、$220~\mathrm{kV}$、$330~\mathrm{kV}$；原稿在 $330~\mathrm{kV}$ 旁注“西北”。
  \item 超高压：$500~\mathrm{kV}$、$700~\mathrm{kV}$；原稿旁注“西北”。
  \item 特高压：$1000~\mathrm{kV}$。
\end{enumerate}

原稿在 $3~\mathrm{kV}$、$6~\mathrm{kV}$ 旁注明“地方用电”，在 $20~\mathrm{kV}$ 旁注明“现代城市”，在 $66~\mathrm{kV}$ 旁注明“东北（用得多）”，并写有“有时写作 $60~\mathrm{kV}$”。

\reviewnote{“均为线电压”应限定在三相系统额定电压的表述中。列表中的三相 $380~\mathrm{V}$ 与单相 $220~\mathrm{V}$ 不能按同一种三相线电压理解；单相用电的电压参考端点需另行说明。}

\reviewnote{原稿超高压第二个数值可辨为 $700~\mathrm{kV}$，常见电网等级则为 $750~\mathrm{kV}$，疑似漏写数字 5，暂保留原稿值，建议核对。国家电网资料使用 $500~\mathrm{kV}$、$750~\mathrm{kV}$ 的表述，参见\href{https://www.nc.sgcc.com.cn/zxzx/zbdt/2024/12/468382.shtml}{国家电网说明}。}

\reviewnote{$60~\mathrm{kV}$ 与 $66~\mathrm{kV}$ 不宜直接当作同一个额定电压的两种写法；需核对教材版本和具体系统。低压至特高压的分类边界以及地域注记，也应以所采用教材和适用标准为准。}

\subsection{规定额定电压的原因}

规定额定电压有利于标准化、技术适配和经济性。

\exampletitle{系统与设备额定电压的一般关系}

% 原图 1-B。
\begin{figure}[H]
  \centering
  \begin{tikzpicture}[every node/.style={font=\small}, notesline]
    \pic at (0.65,0) {notesgenerator};
    \pic at (2.50,0) {notestwowinding};
    \pic at (7.70,0) {notestwowinding};
    \draw (1.03,0) -- (1.96,0);
    \draw (3.04,0) -- (7.16,0);
    \draw (8.24,0) -- (12.15,0);
    \draw[notesflow] (12.15,0) -- (12.15,-0.45);
    \foreach \x in {1.40,3.60,6.60,8.80,11.50}
    \pic at (\x,0) {notesbusbar};
    \node at (0.65,1.25) {$U_{GN}$};
    \node at (1.75,0.80) {$U_{T_1,1N}$};
    \node at (3.45,0.80) {$U_{T_1,2N}$};
    \node at (6.90,0.80) {$U_{T_2,1N}$};
    \node at (8.65,0.80) {$U_{T_2,2N}$};
    \node at (12.15,0.80) {$U$};
    \node at (0.65,-0.80) {发电机};
    \node at (2.50,-0.80) {$T_1$（升压）};
    \node at (5.10,-0.80) {$L_1$};
    \node at (7.70,-0.80) {$T_2$（降压）};
    \node at (10.20,-0.80) {$L_2$};
    \node at (12.15,-0.80) {等值负荷};
    \node[draw, fill=gray!30, minimum width=2.2cm, minimum height=0.65cm]
    at (0.95,-1.65) {系统 1：$U_{N1}$};
    \node[draw, fill=gray!30, minimum width=3.6cm, minimum height=0.65cm]
    at (5.10,-1.65) {系统 2：$U_{N2}$};
    \node[draw, fill=gray!30, minimum width=3.6cm, minimum height=0.65cm]
    at (10.50,-1.65) {系统 3：$U_{N3}$};
    \node at (10.50,-2.35) {负荷电压偏差约 $\pm5\%$（原稿条件）};
  \end{tikzpicture}
  \caption{不同电压等级电网与设备端口电压}
\end{figure}

\reviewnote{图中三个 $U_N$ 表示各段系统额定电压，带设备下标的量表示相应设备或绕组的额定电压；负荷端的 $U$ 则为实际电压。示意图未单独画出厂用电或地方负载支路，不表示发电机侧没有这些负荷。各端口额定值的选择关系仍采用下文的简化条件。}

系统额定电压可直接按各段电网选择，分别记为 $U_{N1}$、$U_{N2}$、$U_{N3}$。设备额定电压需要考虑负荷的电压要求和线路电压降，各段设备的选择不完全相同。

原图在末端负荷旁标注电压范围。以 $U$ 表示该处实际电压，原稿关系整理为
\begin{equation}
  U=(0.95\sim1.05)U_{N3}.
\end{equation}
原稿在设备额定电压的选择依据中，将线路电压降取为约 $5\%$。

\reviewnote{此处约 $\pm5\%$ 的负荷电压范围和约 $5\%$ 的线路电压降，是原稿用于说明额定电压选择的简化条件，不代表所有设备和电网的统一允许偏差或固定压降。}

\examplesubheading{发电机}

考虑厂用电或地方负载的供电要求，原稿将发电机额定电压取为
\begin{equation}
  U_{GN}=1.05U_{N1}.
\end{equation}
原稿注明，即使示意图未画出厂用电或地方负载，也可能存在这些用电部分。

\examplesubheading{升压变压器}

升压变压器 $T_1$ 与发电机直接相接，因此其一次侧额定电压与发电机额定电压相同：
\begin{equation}
  U_{T_1,1N}=U_{GN}=1.05U_{N1}.
\end{equation}
原稿先将二次侧额定电压取为高于相应系统额定电压 $5\%$，再考虑变压器自身约 $5\%$ 的电压降，最终记为
\begin{equation}
  U_{T_1,2N}=1.10U_{N2}.
\end{equation}

% 原图 1-C；中央支路保留原稿连接，并明确标注待核。
\begin{figure}[H]
  \centering
  \begin{circuitikz}[every node/.style={font=\small}, line width=0.65pt]
    \ctikzset{bipoles/length=0.85cm}
    \draw (0,0) node[ocirc] {} -- (0.35,0)
    to[european resistor, l=$R_1$] (2.35,0)
    to[L, l=$X_1$] (4.55,0) -- (5.40,0)
    to[L, l=$X_2'$] (7.70,0)
    to[european resistor, l=$R_2'$] (10.00,0)
    -- (10.80,0) node[ocirc] {};
    \draw (0,-2.80) node[ocirc] {} -- (10.80,-2.80) node[ocirc] {};
    \draw (5.40,0) node[circ] {}
    to[L, l_=$X_0$] (5.40,-1.40)
    to[european resistor, l_=$R_0$] (5.40,-2.80) node[circ] {};
    \draw[dashed, color=reviewgray] (4.85,-0.20) rectangle (6.35,-2.55);
    \node[align=center, color=reviewgray] at (7.55,-1.40)
    {原稿中央支路\\连接与参数待核};
    \draw[notesflow] (-0.50,-0.20) -- (-0.50,-2.60);
    \node[left] at (-0.50,-1.40) {$U_1$};
    \draw[notesflow] (11.30,-0.20) -- (11.30,-2.60);
    \node[right] at (11.30,-1.40) {$U_2'$};
    \draw[notesflow] (0.45,1.10) -- (10.35,1.10)
    node[midway, above] {压降约 $5\%$（原稿估计）};
    \node at (2.20,-3.35) {一次侧};
    \node at (8.70,-3.35) {二次侧折算量};
  \end{circuitikz}
  \caption{变压器等效电路原稿示意（中央支路待核）}
\end{figure}

\reviewnote{图中电阻、电抗符号为整理时补拟；中央支路暂保留原稿可辨的串联画法，不视为已经确认的励磁模型。若它表示通常的铁耗电阻与磁化电抗，二者应按并联支路表示；若采用整体励磁阻抗的串联等效形式，参数含义又不同，需核对教材后确定。参见\href{https://www.rose-hulman.edu/class/ee/HTML/NSF/CCLI_EMD_DUE_0088904/PDFS/lecture_notes_10.pdf}{Rose-Hulman 变压器讲义}。}

\reviewnote{本图采用单相等效电路的表达，撇号表示二次侧折算至一次侧的量。端口电压参考方向均从上端指向下端；它们不能未经绕组接线关系换算，就与前文三相系统的线电压额定值直接混用。上方 $5\%$ 箭头只是原稿采用的压降估计，不表示固定变比或实际测得的恒定压降。}

\reviewnote{额定电压与运行端电压应区分。这里的 $5\%$、$10\%$ 是原稿采用的教材简化取值，不是变压器实际压降的恒定值；实际压降与负载及参数有关。两个 $5\%$ 若按比例逐次相乘，增幅为 $10.25\%$，不能称为精确的 $10\%$；本稿保留原稿约定取值。}

\examplesubheading{线路与降压变压器}

线路 $L_1$ 所属系统的额定电压为 $U_{N2}$。降压变压器 $T_2$ 的一次侧接入该电网，原稿记为
\begin{equation}
  U_{T_2,1N}=U_{N2}.
\end{equation}
当二次侧需要经线路向负荷供电时，原稿取
\begin{equation}
  U_{T_2,2N}=1.10U_{N3}.
\end{equation}
若二次侧直接接电动机，则取
\begin{equation}
  U_{T_2,2N}=1.05U_{N3}.
\end{equation}
负荷（例如电动机）的额定电压则按所在系统额定电压选择：
\begin{equation}
  U_{MN}=U_{N3}.
\end{equation}

原稿特别提示：有些图中直接标注发电机额定电压 $U_{GN}$，而不是系统额定电压 $U_N$；此时，与发电机直接相连的变压器一次侧应与 $U_{GN}$ 一致，不能再重复乘以 $1.05$。

% 手稿 3：MsgID=3042999959834189390，文件修改时间 2026-10-07 19:24:15.914。
\subsection{分接头与实际电压调节}

前例采用的是较理想的电压关系，实际运行中不一定满足用电需求。为调节实际运行电压，可利用变压器分接头改变变比；原稿指出通常在高压侧设置，原因涉及调节准确性和电流较小，并区分有载调压与无载调压。

% 原图 1-D。
\begin{figure}[H]
  \centering
  \begin{circuitikz}[every node/.style={font=\small}, line width=0.65pt]
    \ctikzset{bipoles/length=0.65cm}
    \fill[gray!30] (5.65,1.13) rectangle (9.65,1.67);
    \draw (1.20,2.20) node[ocirc] {} -- (2.00,2.20)
    to[L] (2.00,-1.00) -- (1.20,-1.00) node[ocirc] {};
    \draw (3.25,-1.15) -- (3.25,3.20);
    \draw (3.40,-1.15) -- (3.40,3.20);
    \draw (5.00,3.00) to[L] (5.00,2.20)
    to[L] (5.00,1.40) to[L] (5.00,0.60)
    to[L] (5.00,-0.20) to[L] (5.00,-1.00);
    \foreach \y in {3.00,2.20,1.40,0.60,-0.20}
    \draw (5.00,\y) -- (5.90,\y) node[ocirc] {};
    \draw (5.00,-1.00) -- (5.90,-1.00) node[ocirc] {};
    \fill (5.90,1.40) circle[radius=1.5pt];
    \node[align=center] at (0.35,0.60) {低压侧\\$10.5~\mathrm{kV}$};
    \node at (4.95,3.70) {高压侧分接头};
    \node[anchor=west] at (6.15,3.00) {$+5\%$};
    \node[anchor=west] at (6.15,2.20) {$+2.5\%$};
    \node[anchor=west] at (6.15,1.40) {$0\%$（基准）};
    \node[anchor=west] at (6.15,0.60) {$-2.5\%$};
    \node[anchor=west] at (6.15,-0.20) {$-5\%$};
    \node[anchor=west] at (6.15,-1.00) {公共端};
    \node[anchor=west] at (8.10,3.00) {$254.1~\mathrm{kV}$};
    \node[anchor=west] at (8.10,2.20) {$248.05~\mathrm{kV}$};
    \node[anchor=west] at (8.10,1.40) {$242~\mathrm{kV}$};
    \node[anchor=west] at (8.10,0.60) {$235.95~\mathrm{kV}$};
    \node[anchor=west] at (8.10,-0.20) {$229.9~\mathrm{kV}$};
    \node[align=center] at (6.35,-1.85)
    {百分数基准：$242~\mathrm{kV}$\\右列为按该基准计算的分接标称电压};
  \end{circuitikz}
  \caption{变压器高压侧分接位置及其标称电压}
\end{figure}

原稿以 $-2.5\%$ 分接为例写为
\begin{equation}
  U_{\mathrm{tap}}=242~\mathrm{kV}\times97.5\%.
\end{equation}

\reviewnote{分接头百分数在本例中的基准是 $242~\mathrm{kV}$；图中右列数值是按此基准计算的分接标称电压，不是负荷侧的实际电压。绕组与接点位置为示意，不按实际匝数比例绘制。改变分接位置相当于改变有效匝数和额定变比，不能仅凭“正分接/负分接”的名称就断言负荷侧电压升降；应结合所调绕组和固定电压侧判断。原稿“提高……（准确性、电流小）”处的个别字迹仍待核对。}

\exampletitle{多级电网与三绕组变压器的额定电压}

% 原图 1-E。
\begin{figure}[H]
  \centering
  \begin{subfigure}[t]{0.48\textwidth}
    \centering
    \begin{tikzpicture}[every node/.style={font=\small}, notesline]
      \path[use as bounding box] (-0.40,-4.80) rectangle (6.60,1.50);
      \pic at (0.60,0) {notesgenerator};
      \pic at (2.80,0) {notesthreewinding};
      \draw (0.98,0) -- (2.26,0);
      \draw[notesflow] (3.34,0) -- (6.25,0);
      \pic at (1.65,0) {notesbusbar};
      \pic at (5.65,0) {notesbusbar};
      \draw (2.80,-0.68) -- (2.80,-1.66);
      \pic[rotate=90] at (2.80,-2.20) {notestwowinding};
      \draw (2.80,-2.74) -- (2.80,-3.80) -- (3.76,-3.80);
      \pic[rotate=90] at (2.80,-1.20) {notesbusbar};
      \pic[rotate=90] at (2.80,-3.25) {notesbusbar};
      \pic at (3.35,-3.80) {notesbusbar};
      \pic at (4.30,-3.80) {notestwowinding};
      \draw (4.84,-3.80) -- (5.72,-3.80);
      \pic at (5.25,-3.80) {notesbusbar};
      \pic at (6.10,-3.80) {notesmotor};
      \node at (0.60,-0.85) {$13.8~\mathrm{kV}$};
      \node at (2.80,0.75) {$T_1$};
      \node at (5.10,0.60) {$110~\mathrm{kV}$};
      \node[anchor=west] at (3.35,-1.20) {$35~\mathrm{kV}$};
      \node at (1.70,-2.20) {$T_2$};
      \node[anchor=east] at (2.20,-3.25) {$10~\mathrm{kV}$};
      \node at (4.30,-4.45) {$T_3$};
      \node at (5.25,-3.10) {$0.4~\mathrm{kV}$};
    \end{tikzpicture}
    \caption{电网电压等级与已知发电机额定值}
  \end{subfigure}
  \hfill
  \begin{subfigure}[t]{0.48\textwidth}
    \centering
    \begin{tikzpicture}[every node/.style={font=\small}, notesline]
      \path[use as bounding box] (-0.40,-4.80) rectangle (6.60,1.50);
      \pic at (0.60,0) {notesgenerator};
      \pic at (2.80,0) {notesthreewinding};
      \draw (0.98,0) -- (2.26,0);
      \draw[notesflow] (3.34,0) -- (6.25,0);
      \pic at (1.65,0) {notesbusbar};
      \pic at (5.65,0) {notesbusbar};
      \draw (2.80,-0.68) -- (2.80,-1.66);
      \pic[rotate=90] at (2.80,-2.20) {notestwowinding};
      \draw (2.80,-2.74) -- (2.80,-3.80) -- (3.76,-3.80);
      \pic[rotate=90] at (2.80,-1.20) {notesbusbar};
      \pic[rotate=90] at (2.80,-3.25) {notesbusbar};
      \pic at (3.35,-3.80) {notesbusbar};
      \pic at (4.30,-3.80) {notestwowinding};
      \draw (4.84,-3.80) -- (5.72,-3.80);
      \pic at (5.25,-3.80) {notesbusbar};
      \pic at (6.10,-3.80) {notesmotor};
      \node at (0.60,-0.85) {$13.8~\mathrm{kV}$};
      \node at (1.65,0.65) {$13.8~\mathrm{kV}$};
      \node at (2.80,1.20) {$T_1$};
      \node at (5.10,0.65) {$121~\mathrm{kV}$};
      \node[anchor=west] at (3.35,-0.95) {$38.5~\mathrm{kV}$};
      \node[anchor=west] at (3.35,-1.65) {$35~\mathrm{kV}$};
      \node at (1.70,-2.20) {$T_2$};
      \node[anchor=west] at (3.35,-2.60) {$11~\mathrm{kV}$};
      \node at (3.45,-3.10) {$10~\mathrm{kV}$};
      \node at (4.30,-4.45) {$T_3$};
      \node at (5.25,-3.10) {$0.4~\mathrm{kV}$};
      \node[anchor=east] at (6.50,-4.45) {$0.38~\mathrm{kV}$};
    \end{tikzpicture}
    \caption{各设备与绕组的额定电压}
  \end{subfigure}
  \caption{多级电网及三绕组变压器的额定电压实例}
\end{figure}

\reviewnote{两幅图采用相同接线，左图保留原稿给出的电网电压及发电机额定值，右图在对应绕组端口标出设备额定值，不能将同一位置的两类数值混为实际运行电压。$T_1$ 的三个相交圆表示三绕组变压器，三个绕组对应发电机侧、$110~\mathrm{kV}$ 电网和 $35~\mathrm{kV}$ 电网；图中交叠不表示绕组直接电气相连。末端电动机的 $0.38~\mathrm{kV}$ 仍是原稿按简化条件取的近似值。}

\examplesubheading{发电机与直接相连的变压器}

图中发电机处标注的 $13.8~\mathrm{kV}$ 已是发电机额定电压，已经包含原稿所说的相对于系统额定电压的 $5\%$ 提高，不需要再乘一次。$T_1$ 的低压侧与发电机直接相接，因此
\begin{equation}
  U_{GN}=U_{T_1,\mathrm{LV},N}=13.8~\mathrm{kV}.
\end{equation}
根据原图的系统电压等级及设备注记，$T_1$ 的其余两绕组额定电压分别为
\begin{align}
  U_{T_1,\mathrm{HV},N} & =121~\mathrm{kV},  \\
  U_{T_1,\mathrm{MV},N} & =38.5~\mathrm{kV}.
\end{align}
$T_2$ 与 $T_3$ 的绕组额定电压分别为
\begin{align}
  U_{T_2,1N}/U_{T_2,2N} & =35~\mathrm{kV}/11~\mathrm{kV},  \\
  U_{T_3,1N}/U_{T_3,2N} & =10~\mathrm{kV}/0.4~\mathrm{kV}.
\end{align}

\examplesubheading{电动机额定电压}

电动机直接接在 $T_3$ 的二次侧。原稿认为图中 $T_3$ 出口的 $0.4~\mathrm{kV}$ 已包含高于用电设备额定电压 $5\%$ 的部分，因此电动机额定电压取为
\begin{equation}
  U_{MN}=\frac{0.4~\mathrm{kV}}{1.05}
  \approx0.38~\mathrm{kV}=380~\mathrm{V}.
\end{equation}

\reviewnote{这里需区分图中的 $0.4~\mathrm{kV}$ 究竟是系统额定电压，还是变压器二次绕组额定电压。只有按原稿所采用的直接供电及 $5\%$ 简化条件，才能作上述反推；商不能精确写成 $0.38~\mathrm{kV}$，这里保留约等号。图中“三绕组变压器”的三个绕组，也不能简单套用前例“二次侧”的唯一编号。}

\section{电力系统的接线方式}

原稿参考教材 P5，区分地理接线图与电气接线图。

从供电备用角度，可区分无备用接线与有备用接线。无备用接线多为放射式，结构简单、经济，但供电可靠性相对较低；有备用接线的方式较多，各有优劣。

% 手稿 4：MsgID=7375675206180239227，文件修改时间 2026-10-07 19:24:25.624。
从网络结构和运行方式看，还可区分开式与闭式。原稿提示，应结合供电路径判断备用关系，并记有“闭式设计，开式运行”。

\reviewnote{原稿关于“闭式一定有备用、开式不一定……”的文字有划改，旁边关于供电方向或供电路径的提问也不够清晰，暂保留可确认内容。应区分网络的设计结构、开关的运行状态和实际备用供电能力，后续核对教材对开式、闭式及备用的定义。}

\section{课程体系}

\subsection{分析工具与元件模型}

分析工具包括电路分析和物理概念。通过等值电路等方式建立发电机、变压器、线路和负荷的模型。

\subsection{分析层次}

原稿将分析内容概括为以下层次。

\begin{itemize}
  \item 稳态：研究运行状态，涉及频率 $f$、电压 $V$、有功功率 $P$、无功功率 $Q$ 等。
  \item 电磁暂态：分析、计算故障，由稳态分析扩展到暂态和次暂态分析。
  \item 机电暂态：结合机械与电磁过程，研究同步及稳定性。
\end{itemize}

% 原图 1-H；原图 1-F、1-G 属于用户暂时删除的中性点部分，本次不恢复。
\begin{figure}[H]
  \centering
  \begin{tikzpicture}[
      every node/.style={font=\small},
      toolbox/.style={
          align=center, draw, fill=gray!30,
          text width=2.3cm, minimum height=1.2cm
        },
      modelbox/.style={
          align=center, draw, fill=gray!30,
          text width=2.3cm, minimum height=0.72cm
        },
      analysisbox/.style={
          align=center, draw, fill=gray!30,
          text width=4.8cm, minimum height=1.12cm
        }
    ]
    \node[font=\small\bfseries] at (1.30,1.20) {分析工具};
    \node[font=\small\bfseries] at (5.40,1.20) {元件模型};
    \node[font=\small\bfseries] at (10.20,1.20) {分析层次};
    \node[toolbox] (tools) at (1.30,-1.575) {电路分析\\物理概念};
    \node[modelbox] (generator) at (5.40,0) {发电机};
    \node[modelbox] (transformer) at (5.40,-1.05) {变压器};
    \node[modelbox] (line) at (5.40,-2.10) {线路};
    \node[modelbox] (load) at (5.40,-3.15) {负荷};
    \node[analysisbox] (steady) at (10.20,0)
    {稳态\\运行状态：$f$、$V$、$P$、$Q$ 等};
    \node[analysisbox] (electromagnetic) at (10.20,-1.575)
    {电磁暂态\\故障分析与计算：暂态、次暂态};
    \node[analysisbox] (electromechanical) at (10.20,-3.15)
    {机电暂态\\机械与电磁过程：同步、稳定性};
    \draw[notesline] (tools.east) -- (3.20,-1.575);
    \draw[notesline] (3.20,0) -- (3.20,-3.15);
    \node[align=center] at (3.20,0.75) {建立\\等值电路};
    \foreach \name in {generator,transformer,line,load}
    \draw[notesflow] (3.20,0 |- \name.west) -- (\name.west);
    \draw[notesline] (7.10,0) -- (7.10,-3.15);
    \foreach \name in {generator,transformer,line,load}
    \draw[notesline] (\name.east) -- (7.10,0 |- \name.east);
    \foreach \name in {steady,electromagnetic,electromechanical}
    \draw[notesflow] (7.10,0 |- \name.west) -- (\name.west);
  \end{tikzpicture}
  \caption{电力系统分析的工具、元件模型与分析层次}
\end{figure}

\reviewnote{运行状态变量旁另有一个划改符号，疑似角度记号，暂不确定为 $\delta$，故未加入本图。此处原稿使用电压记号 $V$，前文使用 $U$，后续统一记号时再处理。“稳态、暂态、次暂态”在此表示分析层次的扩展，不能直接当作故障后随时间演变的先后顺序。课程体系图也不表示电磁与机电过程完全独立。}

% 原始手写笔记：MsgID 6861819404153374871，文件时间 2026-10-13 15:41:00。
\chapter{电力系统元件建模}

\section{元件建模的任务与思路}

元件建模的目标是建立能够描述元件外部特性的等值电路，并在准确性与实用性之间取得平衡。基本思路是从元件的物理现象出发，找出对端口特性起主要作用的过程，再用适当的电路元件刻画这些过程。

\reviewnote{本节标题和两条主旨均可辨认；原稿中“描述”之后有一处疑似补写词，暂按“刻画本质”整理。后续使用模型时，应同时注明参考方向、参考侧、相量定义和所采用的近似。}

\section{发电机的建模}

\subsection{定子端口内部模型}

同步发电机空载时，励磁电流 $I_f$ 产生转子磁动势 $F_f$，进而建立主磁通 $\Phi_0$，并在定子绕组中感应空载电势 $E_0$。其因果链可写为
\begin{equation}
  I_f \longrightarrow F_f \longrightarrow \Phi_0 \longrightarrow E_0 .
\end{equation}
负载后，定子电流 $I_a$ 还会产生电枢反应磁通 $\Phi_a$，从而改变气隙中的合成磁场。

\begin{figure}[H]
  \centering
  \begin{tikzpicture}[
    every node/.style={font=\small},
    flux/.style={-{Latex[length=1.8mm]},line width=0.65pt},
    winding/.style={line width=0.8pt}
    ]
    \draw[notesline] (0,0) circle (2.15);
    \draw[notesline,fill=gray!30] (-0.63,-1.12) rectangle (0.63,1.12);
    \draw[notesline] (-0.63,0) -- (0.63,0);
    \node[font=\small\bfseries] at (0,0.55) {N};
    \node[font=\small\bfseries] at (0,-0.55) {S};
    \draw[winding] (-0.82,0.62) arc[start angle=115,end angle=245,radius=0.55];
    \draw[winding] (0.82,-0.62) arc[start angle=-65,end angle=65,radius=0.55];
    \draw[flux] (0.15,1.30) .. controls (1.10,1.08) and (1.12,-1.08) .. (0.15,-1.30);
    \draw[flux] (-0.15,-1.30) .. controls (-1.10,-1.08) and (-1.12,1.08) .. (-0.15,1.30);
    \draw[notesline] (-1.60,1.22) arc[start angle=140,end angle=35,radius=1.15];
    \draw[notesflow] (0.15,2.16) arc[start angle=90,end angle=25,radius=0.62];
    \node at (0.90,2.18) {$\omega_0$};
    \node[align=center] at (3.20,0.75) {定子三相绕组};
    \draw[notesflow] (2.62,0.62) -- (1.72,0.38);
    \node[align=center] at (3.10,-0.28) {转子励磁绕组};
    \draw[notesflow] (2.35,-0.33) -- (0.66,-0.25);
    \node at (-3.15,0.40) {$+$};
    \node at (-3.15,-0.40) {$-$};
    \node at (-3.52,0) {$U_f$};
    \draw[notesline] (-3.02,0.36) -- (-2.20,0.36) -- (-2.20,0.78);
    \draw[notesline] (-3.02,-0.36) -- (-2.20,-0.36) -- (-2.20,-0.78);
    \draw[notesflow] (-2.58,0.36) -- (-2.28,0.36);
    \node at (-2.48,0.68) {$I_f$};
    \node at (1.42,-1.57) {$\Phi_0$};
  \end{tikzpicture}
  \caption{同步发电机励磁与主磁通的示意图}
\end{figure}

\reviewnote{原图为课堂示意图，定子槽形、绕组空间分布和磁通路径均未给出精确结构；上图只保留可确认的转子磁极、直流励磁、旋转方向和主磁通关系。若用于电机理论推导，应对照教材核对极性、绕组轴线和正方向。}

\subsection{每相等值电路与相量关系}

将电枢反应电抗 $X_a$、漏抗 $X_\sigma$ 与定子电阻 $R$ 串联，可得每相稳态等值电路。同步电抗为
\begin{equation}
  X_s=X_a+X_\sigma .
\end{equation}

\begin{figure}[H]
  \centering
  \begin{notescircuit}
    \draw (0,0) to[sV,l_={$\dot E_0$}] (0,2)
    to[L,l={$\mathrm{j}X_a$}] (2.2,2)
    to[L,l={$\mathrm{j}X_\sigma$}] (4.4,2)
    to[R,l={$R$}] (6.4,2) -- (7.2,2);
    \draw (0,0) -- (7.2,0);
    \draw[notesflow] (5.65,2.42) -- (6.50,2.42) node[midway,above] {$\dot I_a$};
    \draw[notesline] (7.2,0) -- (7.2,2);
    \node[right] at (7.2,1.52) {$+$};
    \node[right] at (7.2,0.48) {$-$};
    \node[right] at (7.52,1) {$\dot U$};
    \draw[notesline] (1.42,-0.22) -- (1.42,-0.42) -- (4.98,-0.42) -- (4.98,-0.22);
    \node[below] at (3.20,-0.42) {$X_a+X_\sigma=X_s$};
  \end{notescircuit}
  \caption{同步发电机的每相稳态等值电路}
\end{figure}

按图示发电机电流正方向，有
\begin{equation}
  \dot E_0=\dot U+(R+\mathrm{j}X_s)\dot I_a .
\end{equation}
忽略定子电阻时，相量关系近似为
\begin{equation}
  \dot E_0\approx \dot U+\mathrm{j}X_s\dot I_a .
\end{equation}

\begin{figure}[H]
  \centering
  \begin{tikzpicture}[>=Latex,every node/.style={font=\small}]
    \coordinate (O) at (0,0);
    \coordinate (I) at (5.2,0);
    \coordinate (U) at (3.65,1.80);
    \coordinate (R) at (4.75,1.80);
    \coordinate (E) at (4.75,3.48);
    \draw[notesflow] (O) -- (I) node[below right] {$\dot I_a$};
    \draw[notesflow] (O) -- (U) node[above left] {$\dot U$};
    \draw[notesflow] (U) -- (R) node[midway,below] {$R\dot I_a$};
    \draw[notesflow] (R) -- (E) node[midway,right] {$\mathrm{j}X_s\dot I_a$};
    \draw[notesflow] (O) -- (E) node[above right] {$\dot E_0$};
    \draw[notesline] (1.05,0) arc[start angle=0,end angle=26.3,radius=1.05];
    \node at (1.30,0.27) {$\varphi$};
    \draw[notesline] (1.63,0.80) arc[start angle=26.3,end angle=36.2,radius=1.82];
    \node at (1.72,1.23) {$\theta$};
    \draw[notesline] (0.70,0) arc[start angle=0,end angle=36.2,radius=0.70];
    \node at (0.89,0.47) {$\psi$};
  \end{tikzpicture}
  \caption{同步发电机相量关系示意图}
\end{figure}

图中的三个角满足
\begin{equation}
  \psi=\varphi+\theta .
\end{equation}
原稿给出的功角特性写作
\begin{equation}
  P_e=\frac{E_0U}{X_s}\sin\delta .
\end{equation}

\reviewnote{原稿旁注“书 P42 起”。相量图把 $\theta$ 标在 $\dot U$ 与 $\dot E_0$ 之间，而功率式使用 $\delta$；两者是否指同一个功角需要按课程记号统一。上式是在圆筒转子、忽略电阻的每相模型下得到的每相电磁功率；若 $E_0$、$U$ 取相电压，三相总功率还需乘以 $3$。若采用线电压或凸极机模型，表达式会相应改变。}

% 原始手写笔记：MsgID 1196678885205918328，文件时间 2026-10-13 15:41:09。
\section{变压器的建模}

\subsection{双绕组变压器的等值参数}

双绕组变压器的串联支路由电阻 $R_T=R_1+R'_2$ 和漏抗 $X_T=X_1+X'_2$ 构成，主要反映与电流有关的铜耗和漏磁效应；并联支路由电导 $G_T$ 和电纳 $-\mathrm{j}B_T$ 构成，主要反映与电压有关的铁耗和励磁无功。理想变压器用变比 $k_T$ 表示。

\begin{figure}[H]
  \centering
  \begin{notescircuit}
    \draw (0,0) -- (0,3.0);
    \draw (0,3.0) to[R,l={$R_T$}] (2.2,3.0)
    to[L,l={$\mathrm{j}X_T$}] (4.4,3.0) -- (5.0,3.0);
    \draw (0,0) -- (5.0,0);
    \draw (0.55,3.0) -- (0.55,2.55);
    \draw (0.55,2.55) to[R,l_={$1/G_T$}] (0.55,0.45) -- (0.55,0);
    \draw (1.55,3.0) -- (1.55,2.55);
    \draw (1.55,2.55) to[L,l_={$-\mathrm{j}/B_T$}] (1.55,0.45) -- (1.55,0);
    \draw[notesline] (5.28,1.55) circle (0.55);
    \draw[notesline] (6.08,1.55) circle (0.55);
    \draw[notesline] (4.85,3.0) -- (5.28,2.10);
    \draw[notesline] (4.85,0) -- (5.28,1.00);
    \draw[notesline] (6.08,2.10) -- (6.50,3.0) -- (7.25,3.0);
    \draw[notesline] (6.08,1.00) -- (6.50,0) -- (7.25,0);
    \node at (5.68,0.62) {$k_T$};
    \draw[notesflow] (-0.15,3.42) -- (0.72,3.42) node[midway,above] {$\dot I_1$};
    \draw[notesflow] (6.48,3.42) -- (7.20,3.42) node[midway,above] {$\dot I_2$};
    \node[left] at (0,2.45) {$+$};
    \node[left] at (0,0.55) {$-$};
    \node[left] at (-0.28,1.50) {$\dot U_1$};
    \node[right] at (7.25,2.45) {$+$};
    \node[right] at (7.25,0.55) {$-$};
    \node[right] at (7.52,1.50) {$\dot U_2$};
  \end{notescircuit}
  \caption{双绕组变压器的每相等值电路}
\end{figure}

本节参数均应先折算到同一参考侧。铭牌容量 $S_N$ 为三相容量，铭牌电压 $U_N$ 为线电压；计算采用对称三相的每相 Y 形等值，因此
\begin{equation}
  I_N=\frac{S_N}{\sqrt{3}U_N} .
\end{equation}

\subsection{短路试验}

短路试验通常将一侧短路，从另一侧施加较低电压，使电流达到额定值。由于所加电压较低，可近似忽略励磁支路。短路电压百分数为
\begin{equation}
  U_s\%=\frac{U_s}{U_N}\times100\% .
\end{equation}

\begin{figure}[H]
  \centering
  \begin{notescircuit}
    \draw (0,0) to[sV,l_={$\dot U_s$}] (0,2.4)
    to[R,l={$R_T$}] (2.3,2.4)
    to[L,l={$\mathrm{j}X_T$}] (4.6,2.4) -- (5.3,2.4) -- (5.3,0) -- (0,0);
    \draw[notesflow] (0.65,2.82) -- (1.55,2.82) node[midway,above] {$I_N$};
  \end{notescircuit}
  \caption{变压器短路试验的近似等值电路}
\end{figure}

额定电流下的短路损耗主要是铜耗，故
\begin{align}
  \Delta P_s & =3I_N^2R_T,                     \\
  R_T        & =\frac{\Delta P_sU_N^2}{S_N^2}.
\end{align}
在 $R_T\ll X_T$ 的近似下，有
\begin{align}
  \frac{U_s\%}{100}U_N & \approx\sqrt{3}I_NX_T,                     \\
  X_T                  & \approx\frac{U_s\%}{100}\frac{U_N^2}{S_N}.
\end{align}

\subsection{空载试验}

空载试验通常在低压侧加额定电压，并将另一侧开路。忽略串联支路压降时，测得的空载损耗 $\Delta P_0$ 近似为铁耗。

\begin{figure}[H]
  \centering
  \begin{notescircuit}
    \draw (0,0) to[sV,l_={$\dot U_N$}] (0,2.8) -- (2.8,2.8);
    \draw (0,0) -- (2.8,0);
    \draw (1.05,2.8) to[R,l_={$1/G_T$}] (1.05,0);
    \draw (2.25,2.8) to[L,l_={$-\mathrm{j}/B_T$}] (2.25,0);
    \draw[notesflow] (0.15,3.22) -- (0.88,3.22) node[midway,above] {$\dot I_0$};
  \end{notescircuit}
  \caption{变压器空载试验的近似等值电路}
\end{figure}

空载电流百分数定义为
\begin{equation}
  I_0\%=\frac{I_0}{I_N}\times100\% .
\end{equation}
由有功分量可得
\begin{align}
  \Delta P_0 & =3\left(\frac{U_N}{\sqrt{3}}\right)^2G_T=U_N^2G_T, \\
  G_T        & =\frac{\Delta P_0}{U_N^2}.
\end{align}
若进一步取 $G_T\ll B_T$，则
\begin{align}
  \frac{I_0\%}{100}I_N & \approx\frac{U_N}{\sqrt{3}}B_T,            \\
  B_T                  & \approx\frac{I_0\%}{100}\frac{S_N}{U_N^2}.
\end{align}

变比可由空载时两侧电压之比求得。这里的 $k_T$ 必须与等值电路所采用的电压定义保持一致。

\reviewnote{原稿旁注“书 P23 起”，并把“电压比不是匝数比”作为提醒。更准确的说法是：理想变压器的相电压比等于匝数比；当两侧接线方式不同，线电压比还包含 $\sqrt{3}$ 倍数和相移。上述参数式默认各量已经折算到同一侧且单位自洽；直接把 kW、kV、MVA 或 kVA 的数值混合代入会引入千倍因子，建议先统一成 W、V、VA，或严格沿用教材给出的单位组合。空载电流 $I_0$ 是励磁支路总电流，$B_T$ 的公式还使用了 $G_T\ll B_T$ 的近似。原稿本页小标题字迹不清，现按内容整理为“短路试验”和“空载试验”。}

% 原始手写笔记：MsgID 3944389042640884105，文件时间 2026-10-13 15:41:12。
\subsection{三绕组变压器}

三绕组变压器折算到同一参考侧后，可用三个漏阻抗构成星形等值电路；励磁支路仍接在参考侧端口。

\begin{figure}[H]
  \centering
  \begin{tikzpicture}[every node/.style={font=\small}]
    \coordinate (J) at (4.6,1.6);
    \draw[notesline] (-1.25,1.6) -- (0.8,1.6);
    \draw[notesline,fill=white] (0.8,1.35) rectangle (2.0,1.85);
    \node at (1.4,1.6) {$R_1$};
    \draw[notesline,fill=white] (2.4,1.35) rectangle (3.6,1.85);
    \node at (3.0,1.6) {$\mathrm{j}X_1$};
    \draw[notesline] (2.0,1.6) -- (2.4,1.6) -- (J);
    \draw[notesline] (J) -- (5.2,3.0);
    \draw[notesline,fill=white] (5.2,2.75) rectangle (6.4,3.25);
    \node at (5.8,3.0) {$R_2$};
    \draw[notesline,fill=white] (6.8,2.75) rectangle (8.0,3.25);
    \node at (7.4,3.0) {$\mathrm{j}X_2$};
    \draw[notesline] (6.4,3.0) -- (6.8,3.0) -- (8.35,3.0);
    \draw[notesline] (J) -- (5.2,0.2);
    \draw[notesline,fill=white] (5.2,-0.05) rectangle (6.4,0.45);
    \node at (5.8,0.2) {$R_3$};
    \draw[notesline,fill=white] (6.8,-0.05) rectangle (8.0,0.45);
    \node at (7.4,0.2) {$\mathrm{j}X_3$};
    \draw[notesline] (6.4,0.2) -- (6.8,0.2) -- (8.35,0.2);
    \draw[notesline] (0,1.6) -- (0,0.65);
    \draw[notesline,fill=white] (-0.25,-0.55) rectangle (0.25,0.65);
    \node[left] at (-0.30,0.05) {$1/G_T$};
    \draw[notesline] (0,-0.55) -- (0,-1.0);
    \draw[notesline] (-1.25,1.6) -- (-1.25,0.65);
    \draw[notesline,fill=white] (-1.50,-0.55) rectangle (-1.00,0.65);
    \node[left] at (-1.55,0.05) {$-\mathrm{j}/B_T$};
    \draw[notesline] (-1.25,-0.55) -- (-1.25,-1.0) -- (0,-1.0);
    \node[above] at (-0.62,1.65) {绕组 1};
    \node[right] at (8.35,3.0) {绕组 2，$k_{12}$};
    \node[right] at (8.35,0.2) {绕组 3，$k_{13}$};
    \fill (J) circle (1.3pt);
  \end{tikzpicture}
  \caption{三绕组变压器折算到同一侧的星形等值模型}
\end{figure}

对三个绕组两两进行短路试验，并把测量结果折算到同一容量和电压基准。若三次试验所得成对短路损耗分别为 $\Delta P_{s,12}$、$\Delta P_{s,23}$ 和 $\Delta P_{s,31}$，则
\begin{align}
  \Delta P_{s,12} & =\Delta P_{s,1}+\Delta P_{s,2}, \\
  \Delta P_{s,23} & =\Delta P_{s,2}+\Delta P_{s,3}, \\
  \Delta P_{s,31} & =\Delta P_{s,3}+\Delta P_{s,1}.
\end{align}
因此各绕组的等效短路损耗为
\begin{align}
  \Delta P_{s,1} & =\frac{1}{2}
  \left(\Delta P_{s,12}+\Delta P_{s,31}-\Delta P_{s,23}\right), \\
  \Delta P_{s,2} & =\frac{1}{2}
  \left(\Delta P_{s,12}+\Delta P_{s,23}-\Delta P_{s,31}\right), \\
  \Delta P_{s,3} & =\frac{1}{2}
  \left(\Delta P_{s,23}+\Delta P_{s,31}-\Delta P_{s,12}\right).
\end{align}
在共同基准 $S_B$、参考侧电压 $U_B$ 下，相应电阻可写为
\begin{equation}
  R_i=\frac{\Delta P_{s,i}U_B^2}{S_B^2},
  \qquad i=1,2,3.
\end{equation}

同理，若三组短路电压百分数已换算到同一容量基准，则
\begin{align}
  U_{s,1}\% & =\frac{1}{2}
  \left(U_{s,12}\%+U_{s,31}\%-U_{s,23}\%\right), \\
  U_{s,2}\% & =\frac{1}{2}
  \left(U_{s,12}\%+U_{s,23}\%-U_{s,31}\%\right), \\
  U_{s,3}\% & =\frac{1}{2}
  \left(U_{s,23}\%+U_{s,31}\%-U_{s,12}\%\right).
\end{align}
忽略电阻对短路电压的影响时，漏抗近似为
\begin{equation}
  X_i\approx\frac{U_{s,i}\%}{100}\frac{U_B^2}{S_B},
  \qquad i=1,2,3.
\end{equation}

当三个绕组的额定容量不同，短路试验可能受较小容量绕组的额定电流限制。原稿举例为
\begin{equation}
  S_1=100,\qquad S_2=50,\qquad S_3=100,
\end{equation}
并取 $S_B=100$。若带撇号的损耗是在受限试验容量下测得的值，则损耗按电流平方折算，例如
\begin{align}
  \Delta P_{s,12} & =\left(\frac{100}{50}\right)^2\Delta P'_{s,12},  \\
  \Delta P_{s,23} & =\left(\frac{100}{50}\right)^2\Delta P'_{s,23},  \\
  \Delta P_{s,31} & =\left(\frac{100}{100}\right)^2\Delta P'_{s,31}.
\end{align}

\reviewnote{三绕组参数必须先统一容量基准，再做“和差除以二”。短路损耗随试验电流的平方换算，短路电压百分数则随试验电流线性换算；因此原稿“短路电压不用折算”的旁注不能脱离厂家给出的基准直接使用。原稿在 $R_i$ 式旁的千倍因子有划改，$X_i$ 后的单位也疑似误写为导纳单位，现未保留这些不自洽记号。例中容量未写单位，应结合题目铭牌补全。自耦变压器还涉及公共绕组与串联绕组容量，原稿仅提示“看例 2.5”，本稿暂不补写未经核对的折算式。}

% 原始手写笔记：MsgID 3526514577150184672，文件时间 2026-10-13 15:41:20。
\section{变压器的简化等值模型}

电力系统稳态计算以节点电压和功率为主要变量。由于正常运行时电压变化通常不大，励磁支路消耗的空载有功功率 $\Delta P_0$ 和无功功率 $\Delta Q_0$ 可近似视为常量，并用端口处的恒功率负荷代替励磁导纳。

\begin{figure}[H]
  \centering
  \begin{tikzpicture}[every node/.style={font=\small}]
    \node[align=center] at (-1.55,1.25) {励磁导纳模型};
    \draw[notesline] (0,2.20) -- (1.0,2.20);
    \draw[notesline,fill=white] (1.0,1.95) rectangle (2.25,2.45);
    \node at (1.625,2.20) {$Z_T$};
    \draw[notesline] (2.25,2.20) -- (3.15,2.20);
    \draw[notesline] (0,2.20) -- (0,0.30) -- (3.15,0.30);
    \draw[notesline,fill=white] (-0.25,0.70) rectangle (0.25,1.75);
    \node[left] at (-0.30,1.22) {$Y_T$};
    \draw[notesline] (3.45,1.25) circle (0.52);
    \draw[notesline] (4.20,1.25) circle (0.52);
    \draw[notesline] (3.15,2.20) -- (3.45,1.77);
    \draw[notesline] (3.15,0.30) -- (3.45,0.73);
    \draw[notesline] (4.20,1.77) -- (4.50,2.20) -- (5.10,2.20);
    \draw[notesline] (4.20,0.73) -- (4.50,0.30) -- (5.10,0.30);
    \node at (3.82,0.28) {$k_T$};
    \node[align=center] at (-1.55,-1.60) {恒功率近似};
    \draw[notesline] (0,-0.65) -- (1.0,-0.65);
    \draw[notesline,fill=white] (1.0,-0.90) rectangle (2.25,-0.40);
    \node at (1.625,-0.65) {$Z_T$};
    \draw[notesline] (2.25,-0.65) -- (3.15,-0.65);
    \draw[notesline] (0,-0.65) -- (0,-2.55) -- (3.15,-2.55);
    \draw[notesflow] (0,-0.65) -- (0,-1.75);
    \node[left,align=right] at (-0.15,-1.30)
    {$\Delta\dot S_0=\Delta P_0+\mathrm{j}\Delta Q_0$};
    \draw[notesline] (3.45,-1.60) circle (0.52);
    \draw[notesline] (4.20,-1.60) circle (0.52);
    \draw[notesline] (3.15,-0.65) -- (3.45,-1.08);
    \draw[notesline] (3.15,-2.55) -- (3.45,-2.12);
    \draw[notesline] (4.20,-1.08) -- (4.50,-0.65) -- (5.10,-0.65);
    \draw[notesline] (4.20,-2.12) -- (4.50,-2.55) -- (5.10,-2.55);
    \node at (3.82,-2.57) {$k_T$};
  \end{tikzpicture}
  \caption{用恒功率损耗代替变压器励磁支路}
\end{figure}

该近似可写为
\begin{equation}
  \Delta\dot S_0=\Delta P_0+\mathrm{j}\Delta Q_0\approx\text{常量}.
\end{equation}

\reviewnote{原稿本节第一条的个别字无法确认，现根据后文只整理为“降低模型复杂度”。把励磁支路改为恒功率损耗依赖于运行电压接近额定值；电压偏移较大、研究励磁特性或进行精细损耗计算时，应恢复电压相关的励磁导纳模型。}

\section{变压器的建模（二）：\texorpdfstring{$\Pi$}{Pi} 型电路}

设串联阻抗 $Z_T$ 位于一次侧，把理想变压器一次侧的内部电压记为 $\dot U'_2$，并定义 $\dot U'_2=k_T\dot U_2$。采用图示电流参考方向时，端口方程为
\begin{align}
  \dot U_1 & =Z_T\dot I_1+k_T\dot U_2, \\
  \dot I_1 & =\frac{\dot I_2}{k_T}.
\end{align}
相应传输矩阵为
\begin{equation}
  \begin{bmatrix}
    \dot U_1 \\[2pt]
    \dot I_1
  \end{bmatrix}
  =
  \begin{bmatrix}
    k_T & Z_T/k_T \\[2pt]
    0   & 1/k_T
  \end{bmatrix}
  \begin{bmatrix}
    \dot U_2 \\[2pt]
    \dot I_2
  \end{bmatrix}.
\end{equation}

将端口电流改写为节点电压的函数，有
\begin{align}
  \dot I_1
   & =\frac{k_T}{Z_T}(\dot U_1-\dot U_2)
  +\frac{1-k_T}{Z_T}\dot U_1,            \\
  \dot I_2
   & =\frac{k_T}{Z_T}(\dot U_1-\dot U_2)
  -\frac{k_T(k_T-1)}{Z_T}\dot U_2.
\end{align}
由此得到无磁耦合的 $\Pi$ 型等值电路，其三个支路阻抗为
\begin{align}
  Z   & =\frac{Z_T}{k_T},        \\
  Z_1 & =\frac{Z_T}{1-k_T},      \\
  Z_2 & =\frac{Z_T}{k_T(k_T-1)}.
\end{align}

\begin{figure}[H]
  \centering
  \begin{tikzpicture}[every node/.style={font=\small}]
    \draw[notesline] (0,0) -- (0,2.8) -- (1.4,2.8);
    \draw[notesline,fill=white] (1.4,2.52) rectangle (3.0,3.08);
    \node at (2.2,2.8) {$Z$};
    \draw[notesline] (3.0,2.8) -- (4.4,2.8) -- (4.4,0) -- (0,0);
    \draw[notesline,fill=white] (-0.28,0.85) rectangle (0.28,2.05);
    \node[left] at (-0.36,1.45) {$Z_1$};
    \draw[notesline,fill=white] (4.12,0.85) rectangle (4.68,2.05);
    \node[right] at (4.76,1.45) {$Z_2$};
    \draw[notesflow] (-0.25,3.22) -- (0.80,3.22) node[midway,above] {$\dot I_1$};
    \draw[notesflow] (3.55,3.22) -- (4.60,3.22) node[midway,above] {$\dot I_2$};
    \node[left] at (0,2.35) {$+$};
    \node[left] at (0,0.45) {$-$};
    \node[left] at (-0.30,1.40) {$\dot U_1$};
    \node[right] at (4.4,2.35) {$+$};
    \node[right] at (4.4,0.45) {$-$};
    \node[right] at (4.70,1.40) {$\dot U_2$};
  \end{tikzpicture}
  \caption{变压器的 $\Pi$ 型等值电路}
\end{figure}

这种表示消除了理想变压器造成的磁耦合，便于直接组装节点导纳矩阵。当 $k_T=1$ 时，两条并联支路的导纳均为零，即支路开路，只剩串联阻抗 $Z_T$。

\reviewnote{原稿旁注“P29 起”。由前两式直接消元可知，传输矩阵右上角应为 $Z_T/k_T$；原稿该位置疑似写成 $k_TZ_T$，现按端口方程校正。这里假定 $k_T$ 为实数，且 $Z_T$ 位于所声明的参考侧。若模型包含相移变压器，节点导纳式需要使用复变比及其共轭，不能把上述实数公式直接照搬。$Z_1$、$Z_2$ 是代数等值支路，可能出现负值，并不表示真实的被动元件。}

% 原始手写笔记：MsgID 5541810728296544214，文件时间 2026-10-13 15:41:22。
\section{线路的等值模型}

\subsection{线路类型与四个基本参数}

输电线路主要分为架空线路和电缆线路。架空线路常采用钢芯铝绞线，钢芯主要承受机械拉力，外层铝线主要导电；电缆则由导体、绝缘层、屏蔽或护套等部分组成，具体结构随电压等级和用途而变。

线路中的电流产生发热和磁场效应，分别用电阻 $R$ 和电感 $L$ 描述；线路电压产生电场及绝缘介质中的泄漏效应，分别用电容 $C$ 和电导 $G$ 描述。

\begin{figure}[H]
  \centering
  \begin{tikzpicture}[every node/.style={font=\small}]
    \node[font=\small\bfseries] at (-3.0,2.05) {架空导线截面示意};
    \draw[notesline,fill=gray!30] (-3.0,0.65) circle (0.52);
    \foreach \a in {0,60,...,300}
    \draw[notesline] ({-3.0+0.94*cos(\a)},{0.65+0.94*sin(\a)}) circle (0.42);
    \node[align=center] at (-3.0,-1.00) {内层钢芯\\外层铝线};
    \node[font=\small\bfseries] at (1.35,2.05) {电缆截面示意};
    \draw[notesline,fill=gray!30] (1.35,0.65) circle (1.20);
    \draw[notesline,fill=white] (1.35,0.65) circle (0.84);
    \draw[notesline,fill=gray!30] (1.35,0.65) circle (0.36);
    \draw[notesflow] (2.85,1.25) -- (2.18,1.02);
    \node[right,align=left] at (2.90,1.28) {护套或屏蔽层};
    \draw[notesflow] (2.85,0.45) -- (1.84,0.55);
    \node[right] at (2.90,0.45) {绝缘层};
    \draw[notesflow] (2.85,-0.35) -- (1.60,0.35);
    \node[right] at (2.90,-0.35) {导体};
  \end{tikzpicture}
  \caption{架空导线与电缆的结构示意}
\end{figure}

取单位长度参数 $r_0$、$L_0$、$g_0$ 和 $C_0$，一小段线路可以表示为串联电阻、电感和并联电导、电容。

\begin{figure}[H]
  \centering
  \begin{notescircuit}
    \draw (0,0) -- (0,2.8)
    to[R,l={$r_0\,\mathrm{d}x$}] (2.2,2.8)
    to[L,l={$L_0\,\mathrm{d}x$}] (4.4,2.8) -- (5.3,2.8);
    \draw (0,0) -- (5.3,0);
    \draw (4.7,2.8) to[R,l_={$1/(g_0\,\mathrm{d}x)$}] (4.7,0);
    \draw (5.3,2.8) to[C,l_={$C_0\,\mathrm{d}x$}] (5.3,0);
    \draw[notesflow] (0.20,3.22) -- (1.05,3.22) node[midway,above] {$\dot I$};
  \end{notescircuit}
  \caption{线路微元的分布参数模型}
\end{figure}

\subsection{电阻、电抗与电纳}

均匀导线的单位长度电阻可写为
\begin{equation}
  r_0=\frac{\rho}{A},
\end{equation}
其中 $\rho$ 的长度单位应与 $r_0$ 的“单位长度”相匹配，$A$ 为导体截面积。实际交流电阻还受温度、集肤效应和绞线结构等影响。原稿给出的温度修正式为
\begin{equation}
  r_t=r_{20}\left[1+\alpha(t-20)\right].
\end{equation}
原稿记录的温度系数近似值为
\begin{equation}
  \alpha_{\mathrm{Cu}}=0.00381\ {{}^\circ\mathrm{C}}^{-1},
  \qquad
  \alpha_{\mathrm{Al}}=0.0036\ {{}^\circ\mathrm{C}}^{-1}.
\end{equation}

三相架空线经过完整换位后，等值几何平均距离为
\begin{equation}
  D_{\mathrm{eq}}=\sqrt[3]{D_{AB}D_{BC}D_{CA}}.
\end{equation}
对单根实心圆导线，几何平均半径近似为
\begin{equation}
  r'=r\mathrm{e}^{-1/4}\approx0.779r.
\end{equation}
在原稿采用的工频和单位约定下，单位长度感抗写为
\begin{equation}
  x_0=0.1445\log_{10}\frac{D_{\mathrm{eq}}}{r'}
  \quad \Omega/\mathrm{km}.
\end{equation}

\begin{figure}[H]
  \centering
  \begin{tikzpicture}[every node/.style={font=\small}]
    \coordinate (A) at (0,2.8);
    \coordinate (B) at (-3.1,0);
    \coordinate (C) at (3.8,0.35);
    \fill (A) circle (2pt) node[above] {A};
    \fill (B) circle (2pt) node[below left] {B};
    \fill (C) circle (2pt) node[below right] {C};
    \draw[notesline] (A) -- (B) node[midway,left] {$D_{AB}$};
    \draw[notesline] (B) -- (C) node[midway,below] {$D_{BC}$};
    \draw[notesline] (C) -- (A) node[midway,right] {$D_{CA}$};
    \node[align=center] at (6.2,1.45)
    {完整换位后，各相在一个换位循环内\\依次占据三个几何位置};
    \draw[notesflow] (5.15,1.45) -- (3.95,1.45);
  \end{tikzpicture}
  \caption{三相架空线的相间距离与换位思想}
\end{figure}

若完整换位后各相的自感系数相同、互感系数也可用同一平均值 $M$ 表示，则以 a 相为例，瞬时磁链满足
\begin{align}
  \psi_a(t)            & =L_{aa}i_a(t)+M i_b(t)+M i_c(t), \\
  i_a(t)+i_b(t)+i_c(t) & =0,                              \\
  \psi_a(t)            & =(L_{aa}-M)i_a(t).
\end{align}
因此该对称模型下的每相等值电感为
\begin{equation}
  L_a=L_{aa}-M.
\end{equation}

分裂导线通过增大等效半径来降低导线表面的电场强度，从而减轻电晕及无线电干扰，并可降低线路感抗。对换位线路，原稿给出的单位长度电纳近似式为
\begin{equation}
  b_0=2\pi fC_0
  =\frac{7.58}{\log_{10}(D_{\mathrm{eq}}/r_{\mathrm{eq}})}
  \times10^{-6}\quad \mathrm{S/km}.
\end{equation}

\reviewnote{原稿旁注“书 P10 起”。感抗公式分母处字迹疑似为 $r^2$，但结合紧接着的 $r'=0.779r$，这里整理为几何平均半径 $r'$；请对照教材第 10 页附近确认符号。电纳式中的 $r_{\mathrm{eq}}$ 是电容计算采用的等效半径，不能未经核对就与感抗公式中的几何平均半径混用。$0.1445$ 与 $7.58\times10^{-6}$ 均包含频率、对数底和长度单位等约定，换单位或换频率时需要重新推导。原稿的铜、铝温度系数也应按教材规定的材料状态和温度范围校核。换位后的简化只适用于对称平均模型，不能据此断言未换位线路的三个互感完全相同。}

% 原始手写笔记：MsgID 727534396516291447，文件时间 2026-10-13 15:41:27。
\subsection{集中参数近似与分布参数模型}

线路长度较短时，可以忽略对地电容而采用串联阻抗模型；中等长度线路常采用 $\Pi$ 型集中参数模型；更长线路则需要使用分布参数模型，或用足够多的 $\Pi$ 型单元逼近。

\begin{figure}[H]
  \centering
  \begin{tikzpicture}[every node/.style={font=\small}]
    \node[anchor=east] at (-0.4,2.65) {短线路：};
    \draw[notesline] (0,2.65) -- (0.8,2.65);
    \draw[notesline,fill=white] (0.8,2.38) rectangle (3.2,2.92);
    \node at (2.0,2.65) {$Z=z_0\ell$};
    \draw[notesline] (3.2,2.65) -- (4.0,2.65);
    \node[anchor=east] at (-0.4,0.55) {中等长度线路：};
    \draw[notesline] (0,0.55) -- (0,1.65) -- (0.8,1.65);
    \draw[notesline,fill=white] (0.8,1.38) rectangle (3.2,1.92);
    \node at (2.0,1.65) {$Z=z_0\ell$};
    \draw[notesline] (3.2,1.65) -- (4.0,1.65) -- (4.0,0.55) -- (0,0.55);
    \draw[notesline,fill=white] (-0.25,0.75) rectangle (0.25,1.45);
    \node[left] at (-0.34,1.10) {$Y/2$};
    \draw[notesline,fill=white] (3.75,0.75) rectangle (4.25,1.45);
    \node[right] at (4.34,1.10) {$Y/2$};
    \node[right,align=left] at (5.15,1.10)
    {$Y=y_0\ell$\\$z_0=r_0+\mathrm{j}x_0$\\$y_0=g_0+\mathrm{j}b_0$};
  \end{tikzpicture}
  \caption{短线路和中等长度线路的集中参数模型}
\end{figure}

原稿给出的经验范围是：小于约 $100\,\mathrm{km}$ 时采用短线路模型，约 $100\sim300\,\mathrm{km}$ 时采用一个 $\Pi$ 型，中间长度更大时采用多个 $\Pi$ 型串联或分布参数模型。

单位长度串联阻抗和并联导纳为
\begin{align}
  z_0 & =r_0+\mathrm{j}\omega L_0=r_0+\mathrm{j}x_0, \\
  y_0 & =g_0+\mathrm{j}\omega C_0=g_0+\mathrm{j}b_0.
\end{align}
令坐标 $x$ 从受端指向送端，则微元上的电压、电流关系为
\begin{align}
  \mathrm{d}\dot U & =z_0\dot I\,\mathrm{d}x,                    \\
  \mathrm{d}\dot I & =y_0(\dot U+\mathrm{d}\dot U)\,\mathrm{d}x.
\end{align}
忽略二阶微小量后，得到电报方程
\begin{align}
  \frac{\mathrm{d}\dot U}{\mathrm{d}x} & =z_0\dot I, \\
  \frac{\mathrm{d}\dot I}{\mathrm{d}x} & =y_0\dot U.
\end{align}
进一步有
\begin{align}
  \frac{\mathrm{d}^2\dot U}{\mathrm{d}x^2} & =z_0y_0\dot U, \\
  \frac{\mathrm{d}^2\dot I}{\mathrm{d}x^2} & =z_0y_0\dot I.
\end{align}
定义传播常数和波阻抗为
\begin{align}
  \gamma & =\sqrt{z_0y_0}=\alpha+\mathrm{j}\beta, \\
  Z_c    & =\sqrt{\frac{z_0}{y_0}}
  =R_c+\mathrm{j}X_c
  =\lvert Z_c\rvert\mathrm{e}^{\mathrm{j}\theta_c},
\end{align}
其中 $\alpha$ 为衰减常数，$\beta$ 为相位常数。通解为
\begin{align}
  \dot U(x) & =A_1\mathrm{e}^{\gamma x}+A_2\mathrm{e}^{-\gamma x}, \\
  \dot I(x) & =\frac{A_1}{Z_c}\mathrm{e}^{\gamma x}
  -\frac{A_2}{Z_c}\mathrm{e}^{-\gamma x},
\end{align}
其中 $A_1$、$A_2$ 由边界条件决定。也可写成矩阵形式
\begin{equation}
  \begin{bmatrix}
    \dot U(x) \\[2pt]
    \dot I(x)
  \end{bmatrix}
  =
  \begin{bmatrix}
    A_1 \\[2pt]
    A_1/Z_c
  \end{bmatrix}
  \mathrm{e}^{\gamma x}
  +
  \begin{bmatrix}
    A_2 \\[2pt]
    -A_2/Z_c
  \end{bmatrix}
  \mathrm{e}^{-\gamma x}.
\end{equation}

对架空线路，通常有 $g_0\approx0$ 且 $r_0\ll\omega L_0$。在无损近似下
\begin{align}
  z_0    & \approx\mathrm{j}\omega L_0,          \\
  y_0    & \approx\mathrm{j}\omega C_0,          \\
  Z_c    & \approx\sqrt{\frac{L_0}{C_0}},        \\
  \gamma & \approx\mathrm{j}\omega\sqrt{L_0C_0}.
\end{align}
当受端负荷与波阻抗匹配时，线路只存在正向行波，沿线电压幅值在无损模型中保持不变。用三相线电压 $U$ 表示的自然功率为
\begin{equation}
  P_{\mathrm{nat}}=\frac{U^2}{Z_c}.
\end{equation}
在无损、功率因数接近 $1$ 且端电压条件相同的定性比较下，若 $U_1$、$U_2$ 分别表示送端和受端电压，原稿记录为
\begin{align}
  P>P_{\mathrm{nat}} & :\quad U_2<U_1, \\
  P<P_{\mathrm{nat}} & :\quad U_2>U_1.
\end{align}

\reviewnote{原稿在分布参数等值电路旁注“P18 起”。传播常数又写成疑似 $\beta+\mathrm{j}\alpha$，现按通用记号改为 $\gamma=\alpha+\mathrm{j}\beta$；若课程讲义另有定义，应统一全文。原稿中的 $100\,\mathrm{km}$ 和 $300\,\mathrm{km}$ 是经验分界，不是严格定律：电压等级、电缆电容和计算目的都会改变模型选择。$P$ 与 $P_{\mathrm{nat}}$ 对端电压高低的判断只在相应近似和边界条件下成立，不能直接推广到任意有损线路和任意复功率负荷。长线路无损近似、波阻抗匹配和自然功率的条件可与 NPTEL 的长线路课程材料交叉核对：\href{https://archive.nptel.ac.in/content/storage2/courses/108104051/chapter_2/2_8.html}{Long Transmission Line, NPTEL}。}

% 原始手写笔记：MsgID 923610450715154692，文件时间 2026-10-13 15:41:30。
\section{负荷的等值}

稳态计算中，负荷常采用以下两种基本等值方式。
\begin{enumerate}
  \setlength{\itemsep}{-0.25\baselineskip}
  \setlength{\parsep}{0pt}
  \item 恒功率模型：把负荷表示为给定复功率
        \begin{equation}
          \dot S_L=P_L+\mathrm{j}Q_L.
        \end{equation}
        当系统电压和频率偏离额定值不大时，可用该模型近似描述稳态负荷。
  \item 恒阻抗模型：把负荷表示为
        \begin{equation}
          Z_L=R_L+\mathrm{j}X_L.
        \end{equation}
        该模型常用于需要将负荷直接并入网络阻抗或导纳的计算，例如某些短路计算。
\end{enumerate}

实际负荷由多种设备组成，并随时间和运行方式随机变化。其静态特性更一般地可写成
\begin{align}
  P_L & =F_P(U,f), \\
  Q_L & =F_Q(U,f),
\end{align}
动态负荷还需要引入状态变量和相应的动态方程。工程中的综合负荷模型及参数通常需要根据负荷组成和试验、量测数据确定。

\reviewnote{原稿只给出恒功率、恒阻抗以及“实际负荷具有随机性和动态特性”的提纲，未给出 ZIP 模型比例、频率系数或动态方程，因此本稿不自行补入具体参数。原稿页码只写“下册 P”而未填写数字，建议后续查到何仰赞第四版下册的确切页码后补全。恒功率模型的适用条件也应结合潮流计算的电压范围核对。}

% 原始手写笔记：MsgID 9072798899728588932，文件时间 2026-10-14 20:49:23。
\chapter{标幺制与标幺值}

\section{标幺值的基本概念}

某物理量的标幺值，等于它的有名值与同单位基准值之比，即
\begin{equation}
  x^*=\frac{x}{x_B}.
\end{equation}
标幺值没有量纲。采用标幺制后，各物理量的数量级较为集中，便于计算和比较；若各电压等级的基准值选择恰当，还可以简化变压器两侧参数的折算。

\subsection{基准值的选择}

在 $S$、$U$、$I$、$Z$ 四个量中，通常选取基准容量 $S_B$ 和基准电压 $U_B$，再由二者确定基准电流 $I_B$ 和基准阻抗 $Z_B$。

对单相系统，有
\begin{align}
  I_B & =\frac{S_B}{U_B}, \\
  Z_B & =\frac{U_B}{I_B}=\frac{U_B^2}{S_B}.
\end{align}
对对称三相系统，若 $S_B$ 取三相容量、$U_B$ 取线电压，则
\begin{align}
  I_B & =\frac{S_B}{\sqrt{3}U_B}, \\
  Z_B & =\frac{U_B^2}{S_B}.
\end{align}
因此三相系统仍可直接使用 $Z_B=U_B^2/S_B$，不需要在该式中另加 $\sqrt{3}$。

\subsection{常用物理量的标幺值}

电压、电流、复功率和阻抗的标幺值分别为
\begin{align}
  U^* & =\frac{U}{U_B}, \\
  I^* & =\frac{I}{I_B}, \\
  \dot S^* & =\frac{\dot S}{S_B}
  =\frac{P}{S_B}+\mathrm{j}\frac{Q}{S_B}
  =P^*+\mathrm{j}Q^*, \\
  Z^* & =\frac{Z}{Z_B}
  =\frac{R}{Z_B}+\mathrm{j}\frac{X}{Z_B}
  =R^*+\mathrm{j}X^*.
\end{align}

\reviewnote{原稿旁注“书 P32、P33”。“优点”中的第 1、3 条有部分字迹不清，现根据上下文整理为“便于计算比较”和“简化变压器两侧折算”。三相基准公式不依赖负荷采用 Y 形还是 $\Delta$ 形，但前提是 $S_B$ 为三相容量、$U_B$ 为线电压；实际变压器接线所造成的相移和线、相电压关系仍需单独处理。NPTEL 的标幺制材料同样采用全系统统一的三相容量基准，并使不同电压区的电压基准随变压器变比改变：\href{https://archive.nptel.ac.in/content/storage2/courses/108104051/chapter_1/1_20.html}{Per Unit Representation, NPTEL}。}

\subsection{基准值改变时的换算}

若某阻抗在旧基准下的标幺值已知，可先还原为有名值，再除以新基准阻抗：
\begin{equation}
  Z_{\mathrm{old}}^*
  \longrightarrow
  Z=Z_{\mathrm{old}}^*Z_{B,\mathrm{old}}
  \longrightarrow
  Z_{\mathrm{new}}^*=\frac{Z}{Z_{B,\mathrm{new}}}.
\end{equation}
由此得到
\begin{equation}
  Z_{\mathrm{new}}^*
  =Z_{\mathrm{old}}^*
  \frac{S_{B,\mathrm{new}}}{S_{B,\mathrm{old}}}
  \left(\frac{U_{B,\mathrm{old}}}{U_{B,\mathrm{new}}}\right)^2.
\end{equation}

\reviewnote{原稿只画出了“旧基准标幺值 $\rightarrow$ 有名值 $\rightarrow$ 新基准值 $\rightarrow$ 新标幺值”的换算链，未完整写出最后的合并公式；上式由该换算链和 $Z_B=U_B^2/S_B$ 直接整理得到。使用时应确认新、旧电压基准对应同一电压侧。}

\section{多电压等级网络基准值的逐级归算}

多电压等级网络中，通常在全网选取同一个基准容量 $S_B$，并在某一电压区先选定基准电压，再按照相邻变压器的额定变比逐级确定其他电压区的基准电压。

\begin{figure}[H]
  \centering
  \begin{tikzpicture}[every node/.style={font=\small}]
    \pic at (0,0) {notesbusbar};
    \pic at (1.75,0) {notestwowinding};
    \pic at (3.50,0) {notesbusbar};
    \pic at (5.55,0) {notestwowinding};
    \pic at (7.40,0) {notesbusbar};
    \pic at (9.35,0) {notestwowinding};
    \pic at (11.10,0) {notesbusbar};
    \draw[notesline] (0,0) -- (1.21,0);
    \draw[notesline] (2.29,0) -- (3.50,0) -- (5.01,0);
    \draw[notesline] (6.09,0) -- (7.40,0) -- (8.81,0);
    \draw[notesline] (9.89,0) -- (11.10,0) -- (12.00,0);
    \draw[notesflow] (11.10,0) -- (12.00,0);
    \node[above] at (0,0.48) {$U_{B\mathrm{I}}$};
    \node[above] at (3.50,0.48) {$U_{B\mathrm{II}}$};
    \node[above] at (7.40,0.48) {$U_{B\mathrm{III}}$};
    \node[above] at (11.10,0.48) {$U_{B\mathrm{IV}}$};
    \node[above] at (1.75,0.48) {$T_1$};
    \node[above] at (5.55,0.48) {$T_2$};
    \node[above] at (9.35,0.48) {$T_3$};
    \node[below,align=center] at (1.75,-0.48) {$500/121$\\$\mathrm{kV}$};
    \node[below,align=center] at (5.55,-0.48) {$110/38.5$\\$\mathrm{kV}$};
    \node[below,align=center] at (9.35,-0.48) {$35/11$\\$\mathrm{kV}$};
    \node[below] at (0,-1.05) {$500\,\mathrm{kV}$ 电压区};
    \node[below] at (3.50,-1.05) {$110\,\mathrm{kV}$ 电压区};
    \node[below] at (7.40,-1.05) {$35\,\mathrm{kV}$ 电压区};
    \node[below] at (11.10,-1.05) {$10\,\mathrm{kV}$ 电压区};
  \end{tikzpicture}
  \caption{多电压等级网络的基准电压逐级归算示意}
\end{figure}

设某台变压器两侧额定电压为 $U_{TN1}$、$U_{TN2}$，对应电压区的基准电压为 $U_{B1}$、$U_{B2}$，则其标幺变比为
\begin{equation}
  k_T^*
  =\frac{U_{TN1}/U_{B1}}{U_{TN2}/U_{B2}}.
\end{equation}
若希望该变压器在标幺值等值电路中的变比为 $1$，应令
\begin{equation}
  \frac{U_{B1}}{U_{B2}}
  =\frac{U_{TN1}}{U_{TN2}},
  \qquad k_T^*=1.
\end{equation}

在图示例子中，先取
\begin{equation}
  U_{B\mathrm{I}}=500\,\mathrm{kV},
  \qquad
  U_{B\mathrm{II}}=121\,\mathrm{kV}.
\end{equation}
这与 $T_1$ 的额定电压比 $500/121$ 一致。对 $T_2$，令
\begin{equation}
  k_{T2}^*
  =\frac{110/121}{38.5/U_{B\mathrm{III}}}
  =1,
\end{equation}
可得
\begin{equation}
  U_{B\mathrm{III}}=42.35\,\mathrm{kV}.
\end{equation}

% 原始手写笔记：MsgID 2295871068450974394，文件时间 2026-10-14 20:49:15。
再对 $T_3$ 令
\begin{equation}
  k_{T3}^*
  =\frac{35/42.35}{11/U_{B\mathrm{IV}}}
  =1,
\end{equation}
得到
\begin{equation}
  U_{B\mathrm{IV}}=13.31\,\mathrm{kV}.
\end{equation}

\reviewnote{图中四个电压区下方的 $500$、$110$、$35$、$10\,\mathrm{kV}$ 表示系统电压等级，$U_{B\mathrm{I}}\sim U_{B\mathrm{IV}}$ 是计算采用的基准电压，两者不必数值相等。原稿在全网基准容量附近有一处疑似“400 MVA”的字迹，但该数值没有参与本页电压递推，暂不作为确定条件。$42.35\,\mathrm{kV}$ 和 $13.31\,\mathrm{kV}$ 已按图中三台变压器的额定电压比复算。}

\section{逐级归算法的问题与平均额定电压法}

\subsection{环网中的路径不一致}

在辐射形网络中，从已知电压区沿唯一的变压器链逐级归算即可；在环网中，从不同路径归算到同一电压区时，若变压器额定电压比不完全协调，就可能得到不同的基准电压。

\begin{figure}[H]
  \centering
  \begin{tikzpicture}[every node/.style={font=\small}]
    \draw[notesline] (0,0) circle[radius=0.38];
    \node at (0,0) {$\mathrm{G}$};
    \pic at (1.30,0) {notesbusbar};
    \pic at (3.00,1.55) {notestwowinding};
    \pic at (6.00,1.55) {notesbusbar};
    \pic at (8.55,1.55) {notestwowinding};
    \pic at (10.35,0) {notesbusbar};
    \pic at (3.00,-1.55) {notestwowinding};
    \draw[notesline] (0.38,0) -- (1.30,0);
    \draw[notesline] (1.30,0) -- (1.30,1.55) -- (2.46,1.55);
    \draw[notesline] (3.54,1.55) -- (4.65,1.55);
    \draw[notesline,fill=white] (4.65,1.28) rectangle (5.45,1.82);
    \node at (5.05,1.55) {$X$};
    \draw[notesline] (5.45,1.55) -- (6.00,1.55) -- (8.01,1.55);
    \draw[notesline] (9.09,1.55) -- (10.35,1.55) -- (10.35,0);
    \draw[notesline] (1.30,0) -- (1.30,-1.55) -- (2.46,-1.55);
    \draw[notesline] (3.54,-1.55) -- (10.35,-1.55) -- (10.35,0);
    \node[above] at (1.30,1.90) {$U_{B1}$};
    \node[above] at (6.00,1.90) {$U_{B2}$};
    \node[right] at (10.35,0.30) {$U_{B3}$};
    \node[above] at (3.00,1.95) {$T_1$};
    \node[above] at (8.55,1.95) {$T_2$};
    \node[below] at (3.00,-1.95) {$T_3$};
    \node[below] at (3.00,1.05) {$10.5/242\,\mathrm{kV}$};
    \node[below] at (8.55,1.05) {$220/121\,\mathrm{kV}$};
    \node[above] at (3.00,-1.05) {$10.5/121\,\mathrm{kV}$};
  \end{tikzpicture}
  \caption{不同归算路径造成基准电压不一致的环网示例}
\end{figure}

已知
\begin{equation}
  U_{B1}=10.5\,\mathrm{kV},
\end{equation}
沿上方路径依次令 $T_1$、$T_2$ 的标幺变比为 $1$，有
\begin{align}
  k_{T1}^*
  & =\frac{10.5/10.5}{242/U_{B2}}=1,
  & U_{B2} & =242\,\mathrm{kV}, \\
  k_{T2}^*
  & =\frac{220/242}{121/U_{B3}}=1,
  & U_{B3} & =133.1\,\mathrm{kV}.
\end{align}
若改沿下方的 $T_3$ 归算，则
\begin{equation}
  k_{T3}^*
  =\frac{10.5/10.5}{121/U_{B3}}=1
  \quad\Longrightarrow\quad
  U_{B3}=121\,\mathrm{kV}.
\end{equation}
两条路径分别得到 $133.1\,\mathrm{kV}$ 和 $121\,\mathrm{kV}$，说明无法同时通过基准电压选择使三台变压器的标幺变比都等于 $1$。等价地，沿环路的额定变比乘积为
\begin{equation}
  \frac{10.5}{242}
  \frac{220}{121}
  \frac{121}{10.5}
  =\frac{220}{242}\ne1.
\end{equation}

\reviewnote{原图中 $T_1$ 的高压侧额定值可辨为 $242\,\mathrm{kV}$，$T_2$ 为 $220/121\,\mathrm{kV}$，$T_3$ 为 $10.5/121\,\mathrm{kV}$。课堂旁注指出同一电压等级出现 $242$ 与 $220\,\mathrm{kV}$ 两种额定值，这是两条归算路径不一致的直接原因。环路乘积式为根据原图数据补写的校核式。}

\subsection{平均额定电压近似}

短路电流的近似计算中，可把各电压等级的基准电压取为教材给出的平均额定电压。原稿列出的对应关系为
\begin{equation}
  \begin{array}{c|rrrrrr}
    U_N/\mathrm{kV}
      & 10 & 35 & 110 & 220 & 500 & 1000 \\
    \hline
    U_{\mathrm{av}}/\mathrm{kV}
      & 10.5 & 37 & 115 & 230 & 525 & 1050
  \end{array}.
\end{equation}
这些数值可近似理解为在系统额定电压基础上提高约 $5\%$，但应以课程所用教材的表格为准。

\reviewnote{原稿旁注“P37”，并提出“实际上误差有多大”的问题，但没有给出答案。平均额定电压法是面向特定电压等级和近似计算的约定，表中的 $37/35$、$115/110$ 等比例并不严格等于 $1.05$；不应把“统一乘 $1.05$”当作精确公式。}

\subsection{保留非单位标幺变比}

若既希望使用标幺值，又需要保留实际变比差异，则不强求每台变压器的 $k_T^*$ 都等于 $1$，而是按
\begin{equation}
  k_T^*
  =\frac{U_{TN1}/U_{B1}}{U_{TN2}/U_{B2}}
  =\frac{k_T}{k_B},
  \qquad
  k_B=\frac{U_{B1}}{U_{B2}},
\end{equation}
计算并保留其非单位标幺变比。随后可采用前文的变压器 $\Pi$ 型等值电路，将该变比直接纳入节点导纳模型。这种处理更适合需要反映非标准变比和分接头的潮流计算。

\reviewnote{原稿最后一种方法的名称有涂改，无法可靠辨认为“归一法”或其他术语，故本稿按可确认的物理含义命名为“保留非单位标幺变比”。原稿的处理链可辨为“求标幺值 $\rightarrow$ 归算至某侧 $\rightarrow$ 标幺计算”，并注明 $\Pi$ 型等值电路多用于潮流计算。具体采用有名变比还是标幺变比时，必须同时写清基准电压比 $k_B$，避免把 $k_T$ 与 $k_T^*$ 混用。}

\end{document}
